\chapter{Funkcje i rekurencja}

\section{Po co istnieja funkcje}
Funkcja pozwala nazwac pewien fragment logiki i wywolywac go wiele razy. To jedna z
najwazniejszych technik porzadkowania programu. Zamiast kopiowac ten sam zestaw
instrukcji w kilku miejscach, definiujesz go raz i potem tylko wywolujesz.

Na poziomie praktycznym funkcje w MoonChunk przydaja sie m.in. do:

\begin{itemize}
  \item budowania podpisow i etykiet,
  \item obliczen pomocniczych,
  \item przygotowywania wartosci przed wstawieniem do HTML,
  \item enkapsulacji logiki warunkowej,
  \item reuzycia tych samych transformacji w wielu stronach.
\end{itemize}

\section{Klasyczna deklaracja funkcji}
Podstawowa skladnia wygląda tak:

\begin{lstlisting}[style=moonchunk,caption={Klasyczna funkcja w MoonChunk}]
function greet(name: string): string {
  return "Czesc, " + name;
}
\end{lstlisting}

Mamy tu kilka elementow:

\begin{itemize}
  \item slowo \mc{function},
  \item nazwe funkcji, tutaj \mc{greet},
  \item liste parametrow w nawiasach,
  \item opcjonalny typ zwracany po dwukropku,
  \item ciało funkcji w klamrach.
\end{itemize}

\section{Parametry funkcji}
Parametry to dane przekazywane do funkcji przy wywolaniu.

\begin{lstlisting}[style=moonchunk,caption={Funkcja z dwoma parametrami}]
function fullTitle(title: string, suffix: string): string {
  return title + " - " + suffix;
}
\end{lstlisting}

Wywolanie moze wygladac tak:

\begin{lstlisting}[style=moonchunk,caption={Wywolanie funkcji}]
let result: string = fullTitle("MoonChunk", "Docs");
\end{lstlisting}

Funkcja dostaje argumenty przez wartosc. Dla uzytkownika jezyka oznacza to przede
wszystkim tyle, ze przy kazdym wywolaniu pracuje na danych przekazanych do tego
konkretnego wywolania.

\section{Instrukcja \mc{return}}
\mc{return} konczy dzialanie funkcji i oddaje wynik na zewnatrz.

\begin{lstlisting}[style=moonchunk,caption={Zwracanie wyniku z funkcji}]
function square(n: int): int {
  return n * n;
}
\end{lstlisting}

Jesli funkcja ma cos obliczyc, \mc{return} jest jej naturalnym finalem.

\section{Funkcje bez zwracanej wartosci}
Nie kazda funkcja musi cos oddawac. Czasem celem jest tylko efekt uboczny, np. wypisanie
czegos do konsoli przez \mc{print}. Wtedy przydaje sie typ \mc{void}.

\begin{lstlisting}[style=moonchunk,caption={Funkcja typu void}]
function logMessage(text: string): void {
  print(text);
  return;
}
\end{lstlisting}

\section{Funkcje a HTML}
Najczestszy praktyczny wzorzec polega na tym, ze funkcja przygotowuje wartosc,
a \mc{content} tylko ja wyswietla.

\begin{lstlisting}[style=moonchunk,caption={Funkcja przygotowujaca tekst do strony}]
function badge(isPublished: bool): string {
  return isPublished ? "Opublikowany" : "Szkic";
}

page "" {
  let status: bool = true;
  content {
    <p>Status: {badge(status)}</p>
  };
};
\end{lstlisting}

To bardzo czytelne rozwiazanie. Logika jest zamknieta w funkcji, a HTML pozostaje prosty.

\section{Funkcje jako granica zasiegu}
Kazde wywolanie funkcji tworzy osobna przestrzen dla jej lokalnych nazw. Dlatego ten sam
identyfikator moze bezpiecznie istnieć i na zewnatrz, i wewnatrz funkcji.

\begin{lstlisting}[style=moonchunk,caption={Ta sama nazwa poza i wewnatrz funkcji}]
const x: int = 3;

function demo(): int {
  const x: int = 5;
  return x;
}
\end{lstlisting}

To bardzo przydatne, bo pozwala tworzyc funkcje samodzielne i niekolidujace z otoczeniem.

\section{Rekurencja}
Rekurencja oznacza, ze funkcja wywoluje sama siebie. Brzmi tajemniczo, ale idea jest
prosta: problem rozwiazuje sie poprzez coraz mniejsza wersje tego samego problemu.

Klasyczny przyklad:

\begin{lstlisting}[style=moonchunk,caption={Rekurencyjna funkcja factorial}]
function factorial(n: int): int {
  if (n == 0) {
    return 1;
  };
  return n * factorial(n - 1);
}
\end{lstlisting}

Tutaj:

\begin{itemize}
  \item gdy \mc{n == 0}, funkcja zna odpowiedz od razu,
  \item w przeciwnym razie sprowadza problem do \mc{factorial(n - 1)}.
\end{itemize}

\section{Warunek stopu w rekurencji}
To absolutnie kluczowe. Kazda funkcja rekurencyjna musi miec \important{warunek stopu},
który powoduje, ze w pewnym momencie przestaje wywolywac sama siebie.

Bez tego program wpadlby w nieskonczona sekwencje wywolan.

\section{Rekurencja do przechodzenia po kolejnych krokach}
MoonChunk dobrze radzi sobie z funkcjami, ktore wywoluja same siebie i przy okazji
wypisuja stan kolejnych krokow.

\begin{lstlisting}[style=moonchunk,caption={Rekurencja z licznikiem lokalnym}]
function walk(index: int, max: int): int {
  print("recursive index:", index);
  return (index >= max) ? index : walk(index + 1, max);
}
\end{lstlisting}

To dobry przyklad z dwoch powodow:

\begin{itemize}
  \item pokazuje warunek stopu,
  \item pokazuje, ze kazde wywolanie ma swoj lokalny zestaw argumentow i zmiennych.
\end{itemize}

\section{Zapis strzalkowy}
MoonChunk wspiera tez skrocona, bardziej zwarta skladnie funkcji.

\begin{lstlisting}[style=moonchunk,caption={Funkcja w skladni strzalkowej}]
sum(a: int, b: int): int => a + b;
\end{lstlisting}

Mozesz tez zwrocic blok instrukcji:

\begin{lstlisting}[style=moonchunk,caption={Strzalka z blokiem instrukcji}]
formatTitle(title: string): string => {
  return "[DOCS] " + title;
};
\end{lstlisting}

W praktyce skladnia strzalkowa jest dobra dla krótkich funkcji pomocniczych.
Jesli logika rośnie, klasyczne \mc{function} bywa czytelniejsze.

\section{Funkcje jako wartosci}
MoonChunk pozwala traktowac funkcje jak wartosci, co oznacza, ze mozna je np. wypisywac
przez \mc{print}. To przydatne diagnostycznie i edukacyjnie.

\begin{lstlisting}[style=moonchunk,caption={Wypisanie obiektu funkcji i wyniku jej wywolania}]
function greet(name: string): string {
  return "Hello, " + name;
}

print(greet);
print(greet("MoonChunk"));
\end{lstlisting}

Pierwsza linia pokazuje funkcje jako obiekt/byt wykonywalny, druga -- wynik dzialania.

\section{Najczestsze bledy przy funkcjach}
\begin{itemize}
  \item brak \mc{return} tam, gdzie funkcja ma oddac wynik,
  \item zwracanie wartosci niezgodnej z zadeklarowanym typem,
  \item wywolanie czegos, co nie jest funkcja,
  \item redeklaracja nazw parametrow albo lokalnych zmiennych,
  \item brak warunku stopu w rekurencji.
\end{itemize}

\section{Kiedy wydzielic funkcje}
Jesli ten sam fragment logiki pojawia sie w dwoch miejscach albo jesli blok \mc{content}
zaczyna byc zanieczyszczony obliczeniami, to bardzo dobry moment, zeby wydzielic funkcje.
Dobrze nazwana funkcja poprawia czytelnosc bardziej, niz na pierwszy rzut oka mogloby
sie wydawac.
